\section{Carrier-envelope electrothermal coupling}
\label{sec:carrier-envelope}

For inductive power-transfer applications, electromagnetic carrier periods are
typically much shorter than thermal time scales. vNext exploits this separation
without assuming that electrical operating conditions are constant in time.

\subsection{Two-time-scale model}

Let
\[
  \tau=\omega_c t
\]
be the fast carrier phase and $t$ the slow envelope/thermal time. Port signals
are represented as
\[
  i_{\rm phys}(t)
  =
  \Rea\{i(t)e^{+j\omega_c t}\},
\]
where the complex envelope $i(t)$ varies slowly relative to the carrier:
\[
  \frac{\norm{\dot i}}{\omega_c\norm i}
  \ll1
\]
for the narrowband envelope approximation.

At each slow time, the electromagnetic port relation is
\[
  v(t)
  =
  Z(\omega_c;\mathcal S,T(t))\,i(t).
\]

\subsection{Cycle-averaged heat}

With peak phasors, the cycle-averaged local heat is
\[
 q(\mathbf r,t)
 =
 \frac12
 i(t)\adj H_q(\mathbf r;\mathcal S,T(t))i(t).
 \label{eq:cycle-heat}
\]
The integrated absorbed power is
\[
 P_{\rm abs}(t)
 =
 \frac12
 i(t)\adj D_{\rm abs}(\mathcal S,T(t))i(t).
\]
These are the sources supplied to the slow thermal model. No factor of two is
introduced elsewhere.

\subsection{Circuit-controlled envelope}

For a narrowband external circuit,
\[
  F_{\rm circ}
  \big(
    v(t),i(t),u(t),T(t)
  \big)=0
\]
is solved together with
\[
  v=Zi.
\]
Thus current changes caused by coil resistance, detuning, load variation, or
temperature automatically alter the thermal source.

\subsection{Temperature feedback}

At slow time $t_n$, use the reduced temperature state $z_n$ to evaluate
temperature-sensitive constitutive coordinates
\[
  \vartheta_n=\Theta(z_n,\mathcal S).
\]
Update only the electromagnetic quantities that depend materially on
$\vartheta_n$:
\[
 Z_n=Z(\omega_c;\mathcal S,\vartheta_n),
\]
\[
 H_{q,n}
 =
 H_q(\cdot;\mathcal S,\vartheta_n).
\]
Self/pair caches may use interpolation in passive temperature coordinates.

\subsection{Partitioned time step}

A first implementation uses the following predictor--corrector:
\begin{enumerate}
\item predict $z_{n+1}^{(0)}$ from the previous thermal derivative;
\item evaluate temperature-dependent electromagnetic macromodel at the
      predicted thermal coordinates;
\item solve the small circuit/port problem for $i_{n+1}^{(k)}$;
\item compute loss-channel amplitudes from
      Equation~\eqref{eq:cycle-heat};
\item advance the thermal reduced state implicitly/exponentially;
\item iterate steps 2--5 until the electrothermal state increment and energy
      residual pass tolerances.
\end{enumerate}

The carrier Maxwell field is not time-stepped.

\subsection{When the envelope approximation fails}

If excitation contains switching harmonics, modulation comparable to electrical
memory, or a frequency sweep over which $Z$ varies materially, use the
broadband passive electrical state $x_e$:
\[
 \dot x_e=(J_e-R_e)Q_e x_e+B_e u_e.
\]
Instantaneous/cycle-resolved dissipation is obtained from the passive electrical
storage/dissipation structure and mapped to thermal source channels. The
narrowband envelope model is therefore an efficient specialization of the
broader architecture, not a separate code path with different scene semantics.

\subsection{Electrothermal energy audit}

Over one slow time interval $[t_n,t_{n+1}]$, define
\[
 E_{\rm in}
 =
 \int_{t_n}^{t_{n+1}}
 \frac12\Rea\{v\adj i\}\dd t,
\]
\[
 E_{\rm heat}
 =
 \int_{t_n}^{t_{n+1}}
 P_{\rm abs}(t)\dd t.
\]
The electromagnetic channel decomposition checks
\[
 E_{\rm in}
 =
 E_{\rm heat}
 +E_{\rm out}
 +\Delta E_{\rm em}
\]
with $\Delta E_{\rm em}\approx0$ in a quasi-steady narrowband solve.
The thermal step separately checks
\[
 E_{\rm heat}
 =
 \Delta E_T+E_{\rm env}.
\]
These two identities form an end-to-end energy audit from circuit input to
stored/environmental thermal energy.

\subsection{Thermal time-scale adaptivity}

The time integrator adapts according to the physical thermal energy norm and
changes in electrical envelope/constitutive coordinates. A sharp load change
may force a small thermal step even if the thermal state itself is initially
smooth, because $P_{\rm abs}$ changes discontinuously.

\subsection{Steady operation}

At thermal steady state,
\[
  0=-G_T(z_\ast)z_\ast+B_T(z_\ast)p_\ast+b_{\rm env},
\]
while
\[
  p_\ast
  =
  p\big(
    Z(\omega_c,T_\ast),
    i_\ast,
    T_\ast
  \big)
\]
and the circuit equations hold simultaneously. The steady solver is therefore
a coupled reduced nonlinear problem, not a post-processing solve with frozen
room-temperature impedance.
